\subsection{Efficiency Gap: Utility and Limitations}
\label{sec:eg-limitations}

While the efficiency gap serves as our primary partisan fairness metric, its use requires acknowledging both its utility and limitations---especially in light of Supreme Court skepticism and scholarly critiques.

\subsubsection{The Efficiency Gap Metric}

The efficiency gap, developed by Stephanopoulos and McGhee~\cite{stephanopoulos2015partisan}, measures partisan bias by quantifying asymmetric ``wasted votes'' between parties:

\begin{equation}
\text{EG} = \frac{\text{Wasted Votes}_R - \text{Wasted Votes}_D}{\text{Total Votes}}
\end{equation}

Where wasted votes consist of:
\begin{itemize}
    \item \textbf{Losing party}: All votes cast for losing candidates
    \item \textbf{Winning party}: Surplus votes beyond 50\% $+$ 1 threshold
\end{itemize}

Positive EG indicates Republican advantage (Democrats waste more votes); negative EG indicates Democratic advantage (Republicans waste more votes).

\textbf{Intuition:} Cracking (spreading opposition voters across many districts) and packing (concentrating them in few districts) both create wasted votes. EG quantifies whether one party suffers disproportionate vote wastage.

\subsubsection{Supreme Court Skepticism: \emph{Gill v. Whitford} and \emph{Rucho}}

In \emph{Gill v. Whitford} (2018), the Supreme Court declined to rule on efficiency gap's constitutional adequacy, instead dismissing the case on standing grounds. However, Chief Justice Roberts's opinion in \emph{Rucho v. Common Cause} (2019) expressed broader skepticism about judicially manageable partisan gerrymandering standards~\cite{rucho2019}.

\textbf{Roberts's concerns (paraphrased):}
\begin{enumerate}
    \item \textbf{No clear threshold}: What efficiency gap violates the Constitution? 7\%? 10\%? Different states might require different standards.
    \item \textbf{Temporal sensitivity}: EG can vary substantially across elections as voter preferences shift.
    \item \textbf{Incomplete measure}: EG captures packing and cracking but may miss other gerrymandering tactics.
    \item \textbf{Uniform swing assumption}: EG calculations assume uniform vote swings, which may not reflect reality.
\end{enumerate}

\emph{Rucho} ultimately held partisan gerrymandering claims non-justiciable in federal court, removing efficiency gap (and all other metrics) from federal jurisprudence. However, state constitutional provisions remain viable avenues, and efficiency gap continues to feature in state court litigation.

\subsubsection{Scholarly Critiques}

Beyond judicial skepticism, academic literature identifies specific limitations:

\textbf{1. Close election sensitivity (McGhee \& Stephanopoulos, 2020):}
In elections decided by narrow margins, small vote shifts produce large EG changes. Example: A party winning 51-49\% in ten districts and losing 49-51\% in ten others shows near-zero EG. But if 2\% of voters switch parties, EG could shift by 4-5 percentage points despite unchanged district boundaries.

\textbf{2. Cracking vs. packing ambiguity (Nagle, 2015):}
EG aggregates all wasted votes but doesn't distinguish whether high EG results from cracking (many close losses) or packing (few landslide wins). Courts and reformers may care about these mechanisms differentially.

\textbf{3. Proportionality conflation (Chambers \& Miller, 2013):}
EG approximates but does not equal proportionality. A plan can have zero EG but substantial seat-vote disproportionality (winner's bonus). Conversely, perfectly proportional representation can show non-zero EG depending on vote distributions.

\textbf{4. Turnout differentials (McGhee, 2020):}
EG assumes equal turnout across districts, but urban districts (typically Democratic) often have lower turnout than suburban/rural districts (typically Republican). This means Democrats' wasted votes may be less ``real'' than the metric suggests.

\subsubsection{Why Efficiency Gap Remains Useful}

Despite limitations, efficiency gap retains value for several purposes relevant to our analysis:

\textbf{1. Comparative benchmarking (our primary use):}
We do not claim efficiency gap establishes constitutional thresholds or predicts specific election outcomes. Instead, we use EG to compare algorithmic versus enacted plans under identical electoral conditions. Limitations affecting both plan types equally cancel out in comparative analysis.

\textbf{2. Durability testing (Stephanopoulos \& McGhee, 2015):}
The 7-8\% threshold proposed by Stephanopoulos \& McGhee was never a \emph{constitutional standard} but a \emph{durability threshold}---EG above 7\% tends to persist across multiple elections. Our temporal stability analysis (Section~\ref{sec:temporal-stability}) confirms this: EG remains stable from 2016-2020, demonstrating durable geographic patterns.

\textbf{3. State constitutional litigation (post-\emph{Rucho}):}
While \emph{Rucho} eliminated federal justiciability, state courts retain authority under state constitutions. Pennsylvania, North Carolina, Florida, and other states have struck down partisan gerrymanders under state law, sometimes referencing efficiency gap. Our state-specific baselines provide relevant context.

\textbf{4. Reform advocacy (non-judicial):}
Independent redistricting commissions, good-government groups, and voters evaluating ballot initiatives benefit from quantitative fairness metrics. EG provides accessible intuition (``wasted votes'') for non-experts.

\subsubsection{Methodological Safeguards in Our Analysis}

To mitigate efficiency gap limitations, we employ several methodological safeguards:

\textbf{1. Multiple metrics (Section~\ref{sec:proportionality-analysis}):}
We report efficiency gap, mean-median difference, partisan bias at 50\%, and seats-votes elasticity. Convergence across metrics strengthens conclusions.

\textbf{2. Temporal stability testing (Section~\ref{sec:temporal-stability}):}
Analyzing 2016-2020 elections addresses Roberts's temporal sensitivity concern. Stable EG across elections indicates durable bias, not transient swings.

\textbf{3. State-specific baselines (not universal thresholds):}
We \emph{do not} claim 7\% EG universally indicates unconstitutional gerrymandering. Instead, we provide state-specific algorithmic baselines (e.g., Pennsylvania: $-2.8\%$ algorithmic vs $+7.5\%$ enacted). Deviations from state baselines suggest manipulation.

\textbf{4. Caveat regarding proportionality:}
We explicitly distinguish partisan fairness (symmetric treatment of parties) from proportional representation (seats matching votes). EG primarily measures fairness, not proportionality.

\subsubsection{Reframing Efficiency Gap for Algorithmic Redistricting}

Our use of efficiency gap differs from its original legal purpose:

\textbf{Original purpose (Stephanopoulos \& McGhee):} Detect unconstitutional partisan gerrymanders through EG thresholds (e.g., 7-8\%).

\textbf{Our purpose:} Establish empirical baselines for neutral algorithmic redistricting, demonstrating geographic lower bounds on partisan bias.

\textbf{Key insight:} Algorithmic plans produce $-3.2\%$ EG \emph{without accessing partisan data}. This $-3.2\%$ baseline represents unavoidable geographic sorting under compact districting. Enacted plans showing $+5.1\%$ EG deviate by 8.3 percentage points from this neutral baseline, suggesting human manipulation amplifies geographic patterns.

\textbf{Legal application:} State courts evaluating partisan gerrymandering claims under state constitutions can compare enacted plans to algorithmic baselines. A 10-percentage-point gap (e.g., $+7\%$ enacted vs $-3\%$ algorithmic) provides quantitative evidence that the enacted plan's bias exceeds geographic determinism.

\subsubsection{Alternative Metrics and Robustness}

Section~\ref{sec:proportionality-analysis} presents additional partisan fairness metrics:
\begin{itemize}
    \item \textbf{Mean-median difference}: Detects vote distribution asymmetry
    \item \textbf{Partisan bias at 50\%}: Measures seat advantage at vote parity
    \item \textbf{Seats-votes curves}: Visualize responsiveness and symmetry
    \item \textbf{Declination}: Measures whether parties win districts by similar margins
\end{itemize}

These metrics address different gerrymandering pathologies. Our findings are robust: algorithmic plans show lower bias across \emph{all} metrics compared to enacted plans, not just efficiency gap.

\subsubsection{Conclusion: EG as Comparative Tool}

Efficiency gap has limitations as an absolute constitutional standard---a point the Supreme Court emphasized in \emph{Gill} and \emph{Rucho}. However, for \emph{comparative analysis}, EG remains valuable. Our contribution is not proposing EG thresholds but establishing empirical baselines: neutral algorithms produce $-3.2\%$ EG; enacted plans show $+5.1\%$ EG; the 8.3 percentage point difference quantifies manipulation beyond geography.

State courts, redistricting commissions, and reform advocates can use these baselines to evaluate whether specific enacted plans exhibit bias exceeding what geographic sorting alone would produce. This comparative framing sidesteps many critiques of EG as absolute standard while preserving its utility for assessing partisan fairness.
